\section{Desviaci\`on estandar a largo y corto plazo}
Cada intervalo $R - R$ representa el tiempo entre latidos cardíacos consecutivos, y la secuencia de 
estos intervalos puede analizarse como una serie temporal con propiedades dinámicas \cite{Tayel2015}:

\[
R_1, R_2, R_3, \ldots, R_n
\]

donde cada $R_i$ corresponde al tiempo transcurrido entre el latido $i$ y el latido $i + 1$.

A partir de esta serie pueden obtenerse medidas que caracterizan la variabilidad a corto y largo plazo.

\subsection*{SD1 (variabilidad a corto plazo)}

SD1 cuantifica la variabilidad rápida entre intervalos consecutivos, es decir, la variación instantánea 
de un latido al siguiente \cite{RodriguezBenitez2016}. Se define como:

\[
\text{SD1} = \sqrt{\frac{1}{2} \, \sigma_{\Delta}^2}
\tag{1}
\]

donde $\sigma_{\Delta}$ es la desviación estándar de las diferencias sucesivas $R_{n+1} - R_n$.

Un valor alto de SD1 indica mayor variabilidad \textit{beat-to-beat}, asociada principalmente con la actividad 
parasimpática del sistema nervioso autónomo \cite{Almirall1995, Jimenez2020}.

\subsection*{SD2 (variabilidad a largo plazo)}

SD2 refleja la variabilidad global de la serie temporal, capturando fluctuaciones de mayor escala temporal 
\cite{RodriguezBenitez2016}. Se calcula como:

\[
\text{SD2} = \sqrt{2\sigma^2 - \text{SD1}^2}
\tag{2}
\]

donde $\sigma$ es la desviación estándar total de la serie $R_n$.

Un valor alto de SD2 indica mayor variabilidad a largo plazo, relacionada con cambios emocionales, actividad 
física o alteraciones del ritmo cardíaco \cite{Veloza2019}. Por el contrario, valores bajos pueden asociarse 
con predominio simpático \cite{Jimenez2020}.
